\chapter{Transfers, Incentives and Redistribution}\label{ch:b3:transfers-and-incentives}

Every transfer that is paid to people with low incomes must be withdrawn as their
income rises, or it would be paid to everyone. The withdrawal is a tax on earning,
and for the households that face it, it is often the highest tax in the system.
This chapter reads budget sets under real benefit schedules, sets the predicted
responses against the experiments and reforms that measured them, and states the
formula that turns elasticities and social weights into a tax rate.

\begin{openingchart}{Marginal effective tax rate of a single parent with two
children, by annual earnings, United States, 2024: federal income tax, employee
payroll tax, earned income credit and child tax credit (state taxes and benefits
such as food assistance left out).}{Internal Revenue Service: Revenue Procedure
2023-34, 2024 Instructions for Schedule 8812, Topic 751; our computation}
\begin{tikzpicture}
\begin{axis}[omrecord, width=0.72\linewidth, height=0.4\linewidth, xmin=0,
  xmax=70, ymin=-55, ymax=50, xlabel={Annual earnings (thousands of dollars)},
  ylabel={Marginal rate (\%)}]
  \addplot[omDef, thick, mark=none] table[x expr=\thisrow{mid}/1000, y=mtr] {figdata/bachelor-3/18-transfers-and-incentives/out/mtr.csv};
  \draw[guide] (axis cs:0,0) -- (axis cs:70,0);
\end{axis}
\end{tikzpicture}
\end{openingchart}

% figdata: mtr_stats
Up to 2{,}500 dollars of earnings, each extra dollar earns 40 cents of earned income
credit less 7.65 cents of \omterm{def:b3:efficiency-cost-of-taxation:payroll-tax}{payroll tax}: a marginal rate of $-32$~\%. From 2{,}500 to 17{,}400
dollars the refundable child credit adds 15 cents: $-47$~\%. Between about 25{,}000 and 55{,}800 dollars the credit is
withdrawn at 21~\% while payroll and income taxes are paid: the rate reaches 41~\%.
Above 55{,}800 it falls to 20~\%. The highest marginal rate this parent faces is twice
the one she pays at 60{,}000 dollars, and it falls in the middle of the earnings
distribution, not at the top.

\begin{recall}
The Year 1 volume drew budget sets with kinks and defined the implicit marginal
tax rate: the share of an extra unit of earnings lost to taxes and withdrawn
benefits. The grades 10--12 volume defined the Lorenz curve and the Gini
coefficient, and the Year 2 volume social welfare functions. The previous chapter
defined the \omterm{def:b3:efficiency-cost-of-taxation:eti}{elasticity of taxable income} and the revenue-maximising top rate.
\end{recall}

\begin{notation}[Symbols of this chapter]\label{not:b3:transfers-and-incentives:symbols}
$G$ is a guarantee (the transfer at zero income) and $t$ the rate at which it is
withdrawn; $g$ a \omterm{def:b3:transfers-and-incentives:weights}{social welfare weight}; $a$ and $e$ the Pareto parameter and the
\omterm{def:b3:efficiency-cost-of-taxation:eti}{elasticity of taxable income} of the previous chapter.
\end{notation}

\section{Budget sets under a benefit system}\label{sec:b3:transfers-and-incentives:budget}

\begin{definition}[Participation tax rate; poverty trap]\label{def:b3:transfers-and-incentives:ptr}
\index{participation tax rate}\index{poverty trap}
The \emph{participation tax rate} at earnings $E$ is the share of $E$ lost to taxes and
withdrawn benefits when a person moves from not working to earning $E$: $1 - (N(E) - N(0))/E$,
with $N$ net income. A \emph{poverty trap} is a range of earnings over which the
implicit marginal tax rate is so high that working more raises net income little,
or lowers it.
\end{definition}

\begin{measure}[Marginal effective and participation tax rates]
The marginal rate governs how much to work; the participation rate governs whether
to work at all. They can differ widely: the parent of the opening chart faces a
\omterm{def:b3:transfers-and-incentives:ptr}{participation tax rate} of $-40$~\% at 20{,}000 dollars (she keeps 1.40 dollars of net
income per dollar earned, counting the credits) and of $-6$~\% at 40{,}000, while her
marginal rate there is 41~\%. % figdata: mtr_points
Both are computed for a household type from the statutory schedule; adding the
benefits left out here (food assistance, housing aid, state taxes) raises both, and
each household faces its own.
\end{measure}

\begin{omchart}{Net income of a household of three (a parent and two children) in
Brenland, by gross income, under the 2025 schedule and under the reform of the
weekend problem (dalens a year).}{Invented data: Brenland (\code{country}); our
computation}
\begin{tikzpicture}
\begin{axis}[omrecord, width=0.58\linewidth, height=0.4\linewidth, xmin=0, xmax=40,
  ymin=0, ymax=40, xlabel={Gross income (thousands of dalens)},
  ylabel={Net income (thousands of dalens)}, legend pos=north west,
  legend style={font=\footnotesize}]
  \addplot[omDef, thick, mark=none] table[x expr=\thisrow{gross}/1000, y expr=\thisrow{old}/1000] {figdata/bachelor-3/18-transfers-and-incentives/out/brenland_budget.csv};
  \addplot[omThm, thick, dashed, mark=none] table[x expr=\thisrow{gross}/1000, y expr=\thisrow{new}/1000] {figdata/bachelor-3/18-transfers-and-incentives/out/brenland_budget.csv};
  \legend{2025 schedule, reform}
\end{axis}
\end{tikzpicture}
\end{omchart}

% figdata: brenland_stats
Brenland's income support is paid in full below 12{,}000 dalens and not at all above.
A parent of two who earns 11{,}990 has a net income of 22{,}190; at 12{,}000 she has
13{,}200. She must earn 23{,}240 before she is as well off as just below the line ---
2{,}113 hours a year at 11 dalens an hour, more than a full-time year. A cliff is the
extreme \omterm{def:b3:transfers-and-incentives:ptr}{poverty trap}: a marginal rate above 100~\%.

\section{Measured responses: income guarantees}\label{sec:b3:transfers-and-incentives:nit}

\begin{definition}[Negative income tax]\label{def:b3:transfers-and-incentives:nit}
\index{negative income tax}
A \emph{negative income tax} pays every household a guarantee $G$ withdrawn at a
constant rate $t$ of its income, so that net income is $E + \max(G - tE, 0)$.
\end{definition}

\begin{proposition}[Break-even income]\label{prop:b3:transfers-and-incentives:breakeven}
Under a \omterm{def:b3:transfers-and-incentives:nit}{negative income tax} with guarantee $G$ and rate $t$, the transfer falls to zero
at the break-even income $G/t$. Halving $t$ doubles the number of incomes that receive
something.
\end{proposition}

\begin{proof}
The transfer $G - tE$ is zero at $E = G/t$, positive below. Halving $t$ moves the break-even
to $2G/t$.
\end{proof}

The proposition is the dilemma of every guarantee. A generous $G$ withdrawn slowly
reaches far up the income distribution and costs much; withdrawn fast, it is cheap
but imposes a high marginal rate on everyone below $G/t$.

\begin{model}{The negative income tax}
\assumes A guarantee withdrawn at a constant rate; people choose hours, given
their wage, to maximise utility over consumption and leisure.
\predicts The guarantee raises income at zero hours (an income effect) and the
withdrawal lowers the reward to an hour of work (a substitution effect): both
lower hours, and some people stop working.
\tested Four randomised negative-income-tax experiments in the United States came
to similar results: averaged across them, married men cut their working time by
roughly two weeks of full-time work a year, women by three and young people by four.
% ledger: b3ti:robins.nit
In Finland, 2{,}000 randomly chosen unemployed people received a tax-free 560 euros a
month in 2017 and 2018. In the first year, the clean one, their employment did not
change; over November 2017 to October 2018 they worked 6 more days than the control
group, a period that is mostly the second year. % ledger: b3ti:stm.finland
\breaks The experiments were small and short: a programme known to end changes
little in long-run choices, and none could measure what happens when a whole
labour market is covered. In the Finnish trial the second year, and so most of the
6 days, is confounded with stricter conditions for unemployment benefits that came
in at the start of 2018 for both groups.
\end{model}

\section{In-work credits}\label{sec:b3:transfers-and-incentives:in-work}

\begin{definition}[In-work tax credit; phase-in rate; phase-out rate]\label{def:b3:transfers-and-incentives:in-work}
\index{in-work tax credit}\index{phase-in rate}\index{phase-out rate}
An \emph{in-work tax credit} is a transfer paid only to people with earnings, which
rises with earnings at the \emph{phase-in rate} up to a maximum, stays flat, and is
then withdrawn at the \emph{phase-out rate}.
\end{definition}

The United States credit of the opening chart has, for two children, a \omterm{def:b3:transfers-and-incentives:in-work}{phase-in rate}
of 40~\% (6{,}960 dollars over the first 17{,}400 of earnings) and a \omterm{def:b3:transfers-and-incentives:in-work}{phase-out rate} of
21~\% (the 6{,}960 dollars withdrawn between 22{,}720 and 55{,}768). % ledger: b3ti:irs.eitc
A guarantee lowers the reward to the first hours of work; a credit raises it, and
moves the high marginal rates to earnings where people already work.

\begin{record}[The 1987 expansion of the in-work credit]
The tax reform of 1986 expanded the United States earned income credit from 1987.
Between 1984--1986 and 1988--1989 the labour-force participation of single women with
children rose by 1.4 percentage points, from a base of 73.1~\%, relative to single
women without children, whom the credit did not reach. The hours of women already
working did not change. % ledger: b3ti:el.eitc
\end{record}

The record fits the theory of the phase-out range only in part: the predicted fall in
hours among women in the phase-out did not show. People respond to whether to work
more than to how much, and they may not know the marginal rate the credit imposes
--- the salience of the previous chapter.

\begin{proposition}[Saez's participation result]\label{prop:b3:transfers-and-incentives:participation}
If people respond to taxes mainly by choosing whether to work, the optimal transfer
programme resembles an in-work credit, with negative marginal rates at low earnings
and a small guarantee; if they respond mainly by choosing hours, it is a \omterm{def:b3:transfers-and-incentives:nit}{negative
income tax} with a large guarantee and a high withdrawal rate.
\end{proposition}

\noindent The result is stated here, not proved; it is Saez's (2002) solution of the
optimal-tax problem with both margins. % ledger: b3ti:saez.margins
\emph{A two-type example} shows its logic. With two types --- non-workers and low-paid workers earning $w$ --- a guarantee $G$ to
non-workers withdrawn fully at earnings $w$ raises the \omterm{def:b3:transfers-and-incentives:ptr}{participation tax rate} on the
low-paid to $(G + T)/w$, with $T$ their tax; a credit $C$ paid to low-paid workers lowers it to
$(T - C)/w$. When the number who work falls with the \omterm{def:b3:transfers-and-incentives:ptr}{participation tax rate}, the guarantee
shrinks the base that finances it and the credit enlarges it; with hours fixed, the
intensive margin adds no offsetting cost to the credit.
\cref{exo:b3:transfers-and-incentives:11} puts numbers on it.

\section{The optimal income tax: elasticities and weights}\label{sec:b3:transfers-and-incentives:optimal}

\begin{definition}[Social welfare weight]\label{def:b3:transfers-and-incentives:weights}
\index{social welfare weight}
A \emph{social welfare weight} is the value a social welfare function puts on one more
dalen of consumption for a given person, relative to one more dalen of public
revenue.
\end{definition}

\begin{definition}[Optimal top tax rate]\label{def:b3:transfers-and-incentives:top-rate}
\index{optimal top tax rate}
The \emph{optimal top tax rate} is the marginal rate on incomes above a threshold that
maximises social welfare, given the \omterm{def:b3:efficiency-cost-of-taxation:eti}{elasticity of taxable income}, the shape of the
top of the distribution and the weight on top incomes.
\end{definition}

\begin{proposition}[The top-rate formula]\label{prop:b3:transfers-and-incentives:formula}
With weight $g$ on the consumption of top-bracket taxpayers, Pareto parameter $a$ and
elasticity $e$, the optimal top rate is $\tau = (1 - g)/(1 - g + ae)$.
\end{proposition}

\begin{proof}
Raise the top rate by $d\tau$. Per taxpayer, the mechanical revenue gain is $(z_m - \bar z)\,d\tau$, and
the taxpayer loses the same amount, valued at $g$: a net gain $(1 - g)(z_m - \bar z)\,d\tau$. The
behavioural response lowers revenue by $\tau ez_m\,d\tau/(1 - \tau)$, with no first-order welfare effect on
the taxpayer, who was optimising. At the optimum the two balance: $(1 - g)(z_m - \bar z)(1 - \tau) = \tau ez_m$.
Dividing by $z_m$, $(1 - g)(1 - \tau)/a = \tau e$, so $\tau = (1 - g)/(1 - g + ae)$. With $g = 0$ it is the revenue-maximising
rate of the previous chapter.
\end{proof}

The formula separates what the evidence gives from what it does not. The
elasticity $e$ and the tail $a$ are measured; the weight $g$ is a value judgement, and
the formula says how much each judgement implies, not which to hold. With $a = 2$
and $e = 0.25$, a weight of zero gives 67~\% and a weight of 0.5 gives 50~\%; the debate
about the elasticity, from 0.1 to above 0.8 in the United Kingdom's 50~\% episode,
moves the answer more than the weights do.

\section{Redistribution measured}\label{sec:b3:transfers-and-incentives:redistribution}

Redistribution is measured by comparing inequality of market income (before taxes
and transfers) with inequality of disposable income (after). In Brenland's survey
the Gini coefficient of gross income is 0.290 and that of disposable income 0.223:
taxes and transfers lower it by almost a quarter. % figdata: brenland_stats
Across real economies the comparison needs household surveys that record both
incomes on the same basis. The World Bank's Gini coefficients mix disposable income
and consumption by country and cannot make it; the European survey of income and
living conditions can, for social transfers.

\begin{omchart}{Gini coefficients of equivalised household income, 34 European
countries, 2024 (points): before social transfers, with pensions counted as
transfers and with pensions counted as income, and of disposable income. All three
are measured after direct taxes, so the gaps show what social transfers do, not
taxes.}{Eurostat (ilc\_di12b, ilc\_di12c, ilc\_di12); our sorting}
\begin{tikzpicture}
\begin{axis}[omrecord, width=0.62\linewidth, height=0.95\linewidth, xmin=18, xmax=58,
    ymin=0.3, ymax=34.7, y dir=reverse, ytick=data,
    yticklabels from table={figdata/bachelor-3/18-transfers-and-incentives/out/gini_eu.csv}{country},
    y tick label style={font=\scriptsize}, xlabel={Gini (points)}, grid=none, xmajorgrids,
    legend columns=3, legend style={at={(0.5,1.02)}, anchor=south, font=\footnotesize}]
  \addplot[only marks, mark=*, mark size=1.5pt, omDef] table[x=disposable, y=pos]
    {figdata/bachelor-3/18-transfers-and-incentives/out/gini_eu.csv};
  \addplot[only marks, mark=square*, mark size=1.4pt, omThm] table[x=before_ex_pensions, y=pos]
    {figdata/bachelor-3/18-transfers-and-incentives/out/gini_eu.csv};
  \addplot[only marks, mark=triangle*, mark size=1.7pt, black] table[x=before_all, y=pos]
    {figdata/bachelor-3/18-transfers-and-incentives/out/gini_eu.csv};
  \legend{disposable, before transfers other than pensions, before all transfers}
\end{axis}
\end{tikzpicture}
\end{omchart}

% figdata: gini_eu_stats, gini_eu
The median country's disposable-income Gini is 29.8 points. Social transfers other
than pensions lower the Gini by a median of 4.1 points, from 1.3 in Turkey to 8.0
in Belgium; counting pensions as transfers, the median fall is 16.2 points. The
second number is mostly an artefact of where pensions are counted: a retired person
with a public pension has no income before transfers, so systems that pay pensions
publicly look more redistributive than systems that pay them through private saving.
Neither number includes taxes, which the survey deducts before either Gini. And the
measure is static: it holds incomes before transfers fixed, while the responses of
this chapter change them.

\section{Exercises}\label{sec:b3:transfers-and-incentives:exercises}

\begin{exercise}[$\star$]\label{exo:b3:transfers-and-incentives:1}
In the opening chart, where does the marginal rate peak, and what makes it fall above
55{,}800 dollars?
\end{exercise}

\begin{exercise}[$\star$]\label{exo:b3:transfers-and-incentives:2}
With the opening chart's schedule, the parent's net income is 14{,}360 dollars at
earnings of 10{,}000, 36{,}322 at 30{,}000 and 48{,}349 at 50{,}000. Compute her average gain per
dollar earned between each pair of levels.
\end{exercise}

\begin{exercise}[$\star$]\label{exo:b3:transfers-and-incentives:3}
A \omterm{def:b3:transfers-and-incentives:nit}{negative income tax} has a guarantee of 8{,}000 dalens and a rate of 40~\%. Find the
break-even income, and the transfer at earnings of 5{,}000 and 15{,}000.
\end{exercise}

\begin{exercise}[$\star$]\label{exo:b3:transfers-and-incentives:4}
What data on hours would contradict the \omterm{def:b3:transfers-and-incentives:nit}{negative income tax}'s prediction?
\end{exercise}

\begin{exercise}[$\star\star$]\label{exo:b3:transfers-and-incentives:5}
In Brenland, a parent of two who does not work has a net income of 10{,}200 dalens.
Compute her \omterm{def:b3:transfers-and-incentives:ptr}{participation tax rate} at gross incomes of 8{,}000 and 15{,}000 under the 2025
schedule.
\end{exercise}

\begin{exercise}[$\star\star$]\label{exo:b3:transfers-and-incentives:6}
Find the optimal top rate with $g = 0.2$, $a = 2$ and $e = 0.25$. How does it change if $e$ is 0.5?
\end{exercise}

\begin{exercise}[$\star\star$]\label{exo:b3:transfers-and-incentives:7}
The \omterm{def:b3:transfers-and-incentives:in-work}{phase-out rate} of an in-work credit is lowered, with the same maximum. Who is
expected to work more, and who less? What happens to the cost?
\end{exercise}

\begin{exercise}[$\star\star$]\label{exo:b3:transfers-and-incentives:8}
Compute, with the 2024 parameters, the marginal effective tax rate of a single parent
with one child from 0 to 60{,}000 dollars of earnings. Where is it highest, and how
does it compare with the two-child schedule?
\end{exercise}

\begin{exercise}[$\star\star$]\label{exo:b3:transfers-and-incentives:9}
What evidence would contradict the claim that the participation of single mothers
responds strongly to taxes?
\end{exercise}

\begin{exercise}[$\star\star\star$]\label{exo:b3:transfers-and-incentives:10}
Prove the top-rate formula, and show that it is decreasing in $g$, $a$ and $e$. Why does the
behavioural response have no first-order effect on the taxpayer's welfare?
\end{exercise}

\begin{exercise}[$\star\star\star$]\label{exo:b3:transfers-and-incentives:11}
A population of 100 can earn 10{,}000 by working; a share $1 - \mathrm{PTR}$ of them works, where PTR is
the \omterm{def:b3:transfers-and-incentives:ptr}{participation tax rate}. The budget for transfers is 200{,}000. Compare (a) a
guarantee paid to non-workers only, set so as to use the whole budget, with (b) a
credit paid equally to all workers: the number who work and the transfer per
recipient, assuming no tax otherwise.
\end{exercise}

\begin{exercise}[$\star\star\star$]\label{exo:b3:transfers-and-incentives:12}
With Brenland's survey, compute the Gini coefficients of gross and of disposable
income, and the share of the reduction due to taxes and to transfers (compute the
Gini of gross income minus tax).
\end{exercise}

\section{Problem: Reforming Brenland's Income Support}\label{sec:b3:transfers-and-incentives:problem}

\begin{problem}[{Weekend problem --- Reforming Brenland's income support}]\label{pb:b3:transfers-and-incentives:brenland}
Brenland's 2025 income support pays $6{,}000 + 1{,}500(n - 1)$ dalens a year to a household of
$n$ members whose gross income is below 12{,}000, and nothing above. A child benefit of
1{,}200 per member beyond two is paid to all. Income tax follows the bands of the
previous chapter. The household survey (\code{load("households")}) records the
transfers. The wage of a low-paid parent is 11 dalens an hour. % figdata: brenland_stats, brenland_budget

\medskip\noindent\textbf{Part I --- the current schedule.}
\begin{enumerate}
\item Compute the support and the net income of a non-working parent of two
($n = 3$).
\item Compute her net income at gross incomes of 11{,}990 and 12{,}000.
\item What is the implicit marginal tax rate on the last 10 dalens?
\item Below 12{,}000, what is her marginal rate?
\item How many of the survey's 2{,}000 households are below the line?
\item Rebuild the survey's \code{transfers} column from the rules.
\end{enumerate}

\medskip\noindent\textbf{Part II --- the poverty trap.}
\begin{enumerate}[resume]
\item Find the gross income at which she is as well off as at 11{,}990.
\item How many hours a year must she work to reach it?
\item Compute her \omterm{def:b3:transfers-and-incentives:ptr}{participation tax rate} at 15{,}000.
\item Why might she choose to earn 11{,}990 rather than 15{,}000?
\item What would the survey show around 12{,}000 if households respond?
\item Why is the trap worse for larger households?
\end{enumerate}

\medskip\noindent\textbf{Part III --- the reform.} The support is withdrawn at 50~\% of
gross income instead of all at once, and an in-work credit pays 20~\% of the earnings of
households that work, up to 2{,}000, withdrawn at 10~\% of gross income above 20{,}000.
\begin{enumerate}[resume]
\item Find the break-even income of the support for $n = 3$.
\item Compute her net income at 15{,}000 after the reform.
\item Compute her \omterm{def:b3:transfers-and-incentives:ptr}{participation tax rate} at 15{,}000 after the reform.
\item Compute her \omterm{def:b3:transfers-and-incentives:ptr}{participation tax rate} at 8{,}000, before and after.
\item Who loses from the reform? Give an example household.
\item Compute the static cost of the reform per household of the survey.
\item How many households gain and how many lose?
\end{enumerate}

\medskip\noindent\textbf{Part IV --- redistribution.}
\begin{enumerate}[resume]
\item Compute the Gini coefficient of gross income.
\item Compute it for disposable income, before and after the reform.
\item By how much does the reform reduce inequality as measured?
\item Why is this a static measure?
\item What responses would change the cost of the reform, and in which direction?
\item State the result: the \omterm{def:b3:transfers-and-incentives:ptr}{participation tax rate} at 15{,}000 before and after the
reform, and its cost per household.
\end{enumerate}
\end{problem}
